% hardware.tex -- Methodology: hardware and experimental set-up (draft 4 Oct 2026).

\subsection{Hardware and experimental set-up}
\label{sec:hardware}

\paragraph{UAVs.} Four quadrotors run Betaflight in rate mode on six-cell batteries. Each has its own
ExpressLRS receiver and a ground transmitter on a USB serial port, so every UAV is commanded on its
own link at 100\,Hz. The throttle, three rate sticks, the arm switch and an auxiliary channel are
sent; the auxiliary channel switches the UAV's magnet and is held on from start-up, so a quiet link
cannot drop the payload. The flight controller's telemetry (battery voltage, flight mode, link
quality) comes back on the same link. As flown, each UAV weighs 0.55\,kg: an airframe and battery of
0.475\,kg, plus its rod and magnet.

\paragraph{Rods and payload.} Each UAV carries a rigid rod of 0.55\,m and 0.075\,kg with a ball joint
at both ends. The upper joint sits 4\,cm below the UAV's centre. The payload is a 0.86\,kg ring of
0.51\,m diameter and 40\,mm thickness, with twelve magnet plates every 30$^\circ$ on its top face at a
radius of 0.225\,m.
Each rod ends in an electromagnet that holds on a plate with 10--12\,N and releases under a jerk.
A layout is named by its plate numbers, anticlockwise from plate 0 on the ring's $+x$ axis
(Figure~\ref{fig:geometry}).

% fig_geometry.tex -- M3: payload and rod geometry (TikZ). Dimensions from the rig (30 Sep 2026):
% plates at r 0.225 m on a 500 mm ring, 12 plates every 30 deg anticlockwise from +x, rod 0.55 m,
% pivot 4 cm below the UAV centre, working elevation 45 deg. Drawn to scale (1 m = 9 cm).
\begin{figure}[h]
\centering
\begin{tikzpicture}[x=9cm, y=9cm, >=Stealth, font=\small]
  % ---------------- (a) top view ----------------
  \begin{scope}
    \fill[gray!12] (0,0) circle (0.25);
    \fill[white] (0,0) circle (0.19);
    \draw[gray!60, thin] (0,0) circle (0.25) (0,0) circle (0.19);
    % the even triangle left after detaching plate 3
    \draw[densely dashed, gray!70] (30:0.225) -- (150:0.225) -- (270:0.225) -- cycle;
    % plates
    \foreach \k in {0,...,11} {
      \pgfmathsetmacro{\a}{30*\k}
      \node[draw, fill=white, minimum size=0.11cm, inner sep=0pt, rotate=\a] at (\a:0.225) {};
      \node[font=\scriptsize, text=gray!70!black] at (\a:0.29) {\k};
    }
    % the four carrying plates
    \foreach \k in {1,3,5,9} { \fill[black] (30*\k:0.225) circle (0.012); }
    \draw[thick, black] (90:0.225) circle (0.022);
    % centre, axes, rho
    \fill (0,0) circle (0.008);
    \draw[->, thin, gray] (0,0) -- (0.09,0) node[below, font=\scriptsize, text=black] {$x_L$};
    \draw[->, thin, gray] (0,0) -- (0,0.09) node[left, font=\scriptsize, text=black] {$y_L$};
    \draw[->, thick] (0,0) -- (30:0.213) node[pos=0.62, above left, inner sep=1pt] {$\boldsymbol{\rho}_i$};
    \node[font=\scriptsize, align=left, anchor=north] at (0,-0.33) {$\bullet$ carrying plates 1, 3, 5, 9\quad $\bigcirc$ plate 3, detached};
    \node[anchor=north west] at (-0.36,0.36) {(a)};
  \end{scope}
  % ---------------- (b) side view of one rod (A plate, P pivot, C UAV centre) ----------------
  \begin{scope}[shift={(0.62,-0.30)}]
    \fill[gray!12] (-0.06,0) rectangle (0.0,0.03);
    \draw[gray!60, thin] (-0.06,0) rectangle (0.0,0.03);
    \fill[black] (-0.012,0.03) rectangle (0.012,0.038);
    \coordinate (A) at (0,0.034);
    \coordinate (P) at (0.3889,0.4229);
    \coordinate (C) at (0.3889,0.4629);
    \draw[very thick] (A) -- (P);
    \fill (A) circle (0.007) (P) circle (0.007);
    \draw[thick] (0.2989,0.4629) -- (0.4789,0.4629);
    \draw[gray] (0.2989,0.4689) ellipse (0.035 and 0.006) (0.4789,0.4689) ellipse (0.035 and 0.006);
    \fill (C) circle (0.006);
    \draw[->, thick] (A) -- (0.1131,0.1471) node[pos=0.55, above left, inner sep=1pt] {$\mathbf{q}_i$};
    \draw[gray] (0.16,0.034) arc (0:45:0.16);
    \draw[gray, thin] (A) -- (0.22,0.034);
    \node[font=\scriptsize] at (0.185,0.1105) {45$^\circ$};
    \node[font=\scriptsize, anchor=north west] at (0.2045,0.2185) {$l_i = 0.55$\,m};
    \draw[gray, thin, densely dotted] (P) -- (0.5689,0.4229) (C) -- (0.5689,0.4629);
    \draw[<->, thin] (0.5689,0.4229) -- (0.5689,0.4629) node[midway, right, font=\scriptsize] {$\boldsymbol{\delta}$: 4\,cm};
    \node[font=\scriptsize, anchor=north] at (-0.03,-0.005) {ring};
    \node[font=\scriptsize, anchor=south] at (0.3889,0.49) {UAV $i$};
    \node[anchor=north west] at (-0.08,0.66) {(b)};
  \end{scope}
\end{tikzpicture}
\caption[Payload and rod geometry]{Payload and rod geometry, to scale. (a) Top view of the 0.86\,kg ring
with its twelve magnet plates at a radius of 0.225\,m, numbered anticlockwise from the $x_L$ axis. The
four-UAV layout uses plates 1, 3, 5 and 9; detaching the UAV on plate 3 leaves the even triangle
(dashed), under which the ring hangs level. (b) Side view of one rod at the working elevation: the rod
direction $\mathbf{q}_i$ points from the plate to the UAV, and the rod pivots $\boldsymbol{\delta}$ = 4\,cm
below the UAV's centre.}
\label{fig:geometry}
\end{figure}


\paragraph{Motion capture.} OptiTrack tracks one rigid body per UAV and one for the ring. Its poses
are the only state measurement: the velocities and body rates of the UAVs and of the ring are
differentiated from them.

\paragraph{Thrust map.} The trackers need the throttle that produces a given thrust. A proportional
map with one gain, identified at hover, did not hold across loads and battery states. The map was
therefore identified on the rig: each UAV hovered free with hung masses of 0, 98 and 196\,g over
battery voltages of 22.7--24.6\,V. Thirteen points fit
\begin{equation}
    u_T = a_i + 0.507\,M - 0.022\,(V - 23.5),
\end{equation}
with $u_T$ the throttle, $M$ the mass carried (kg), $V$ the battery voltage and $a_i$ = 0.181--0.187
for the four UAVs, with an rms residual of 0.003. Above its offset the map has a gain of 35.2\,m/s$^2$
per unit throttle for a 0.55\,kg UAV, which is the $k_T$ of the tracker model \eqref{eq:trk_model}.
With this map and one shared offset, three of the four UAVs hovered free within 1\,cm of their targets
with no integral action, and the fourth, whose own offset is the lowest, 3.7\,cm high. A map
identified at hover with one proportional gain had left the UAVs 14--17\,cm low.

Table~\ref{tab:params} lists the parameters the rig flies with.

\begin{table}[htbp]
\centering
\caption{Rig parameters, from the rig's launch profile. The rig-twin simulations use the same values.}
\label{tab:params}
% generated by tools/thesis_sim_tables.py
\begin{tabular}{llc}
\toprule
Parameter & Value & Unit \\
\midrule
\multicolumn{3}{l}{\textit{Fleet and payload}} \\
\quad UAV mass, with rod and magnet & 0.55 & kg \\
\quad payload mass & 0.86 & kg \\
\quad rod length & 0.55 & m \\
\quad rod pivot below the UAV centre & 0.04 & m \\
\quad attachment radius & 0.225 & m \\
\multicolumn{3}{l}{\textit{Thrust map}} \\
\quad throttle offset $a_i$ (shared) & 0.185 &  \\
\quad voltage slope & 0.022 & 1/V \\
\quad gain above the offset, $k_T$ & 35.2 & m/s$^2$ \\
\quad throttle cap & 0.8 &  \\
\multicolumn{3}{l}{\textit{Takeoff and lift}} \\
\quad creep speed & 0.2 & m/s \\
\quad rod elevation at the hand-over & 45 & \si{\degree} \\
\quad pretension & 3 & s \\
\quad tautness gate, measured over typed rod & 0.9 &  \\
\quad lift ramp & 0.1 & m/s \\
\multicolumn{3}{l}{\textit{Height integral}} \\
\quad gain $k_i$ & 0.4 & 1/s \\
\quad bound & 0.3 & m \\
\quad gate: height error under & 0.25 & m \\
\multicolumn{3}{l}{\textit{Landing}} \\
\quad descent speed & 0.2 & m/s \\
\quad touchdown tolerance & 0.1 & m \\
\multicolumn{3}{l}{\textit{Reconfiguration}} \\
\quad fewest UAVs after a detach & 3 &  \\
\quad weld radius & 0.08 & m \\
\multicolumn{3}{l}{\textit{Watchdogs}} \\
\quad motion-capture timeout & 0.25 & s \\
\quad planner solve budget & 0.06 & s \\
\bottomrule
\end{tabular}

\end{table}

\paragraph{Arming.} Betaflight refuses to arm while any of its arming checks fails, and then waits
for the arm switch to go low and high again. A switch raised at such a moment is ignored until it is
toggled, which is the likely reason two of the UAVs often failed to arm on the first attempt on the
rig and armed after a power cycle. Each UAV's radio node therefore holds the arm switch low for
0.3\,s before raising it, confirms the arm from the flight mode the flight controller reports, and
retries with a fresh switch edge up to three times. This has been tested against a simulated flight
controller, not yet on the rig. The fleet manager waits
for every flight controller to report armed and names any UAV that did not, with the reason it
reported, before TAKEOFF is accepted.

\paragraph{Preflight checks.} Before each flight a script checks, from the live topics:
motion capture for every UAV and the payload (fresh, above 20\,Hz, no missing values); each UAV
resting level (under 20$^\circ$) and the payload resting level (under 25$^\circ$); every battery
above 23.6\,V and its telemetry present; the same magnet channel on every radio; the parameters the
controllers actually loaded (masses, rod length, attachment radius, thrust map, throttle cap);
each UAV within tolerance of its typed attachment azimuth; and each rod length from motion capture.
It reports GO or NO-GO.

\paragraph{Flight procedure.} Each flight is started with one command that clears leftover
processes, records every relevant topic and writes the logs to a per-flight folder, so that every
flight can be read back before the next. The operator holds a disarm button on the ground station;
a disarm or a fault on any UAV disarms the whole fleet.

\paragraph{Test ladder.} New behaviour is flown in an order that isolates one effect at a time:
free hover, to check the thrust map; then the payload held at its lift height; hover; the circle and
the figure-8; the detach; and the attach last. Each step is first flown in the rig twin with the
same launch commands.
